\section{Localized modes and the finite fermion model}\label{sec:wells}

The positive spin model now fixes the target exchange strengths $K_e$.
This section constructs smooth localized modes and a finite fermion
Hamiltonian whose low-energy spin interaction reproduces $H_+$ up to
scale and a scalar shift. Its occupancy
penalty retains the intersite direct Coulomb coefficients. Computing the
cost of virtual double occupancy determines the required hopping, which
we then obtain by choosing the mode separations.

We work in coordinates \(y=\lambda x=(r,z)\) with \(r\in\mathbb R^2\). The common \(\lambda\) will be polynomially large. After dividing energies by \(\lambda^2\), kinetic energy is the usual \(-\Delta/2\), each unit nucleus has coupling \(1/\lambda\), and electron repulsion has coefficient \(1/\lambda\).

\subsection{Local modes and Coulomb coefficients}\label{subsec:local-modes}

Take \(f(r)=(1-|r|^2)_+^{16}\), and fix
\[
 \phi=(-\Delta_r+1)^{-1}f,\qquad \varphi=\phi/\|\phi\|_2,\qquad
 W(r)=-\frac{f(r)}{2\phi(r)} .
\]
These functions have more bounded continuous derivatives than needed below (in particular through order six). \(W\) is nonpositive and compactly supported, and \(\phi\) is strictly positive, radial, exponentially decaying. For details,
\[
  \phi(r)=\int_0^\infty e^{-t}\int_{\mathbb R^2}
   (4\pi t)^{-1}e^{-|r-b|^2/(4t)} f(b)\,db\,dt
\]
solves the equation by the Gaussian heat kernel identity. For \(|r|>2\), it and the derivatives needed below are bounded in absolute value by \(C(1+|r|)^C e^{-|r|}\), by differentiating on \(f\) and using \(t+(|r|-1)^2/(4t)\ge |r|-1\); for example integrate \(dt/t\) with this bound from 1 to \(|r|^2\), bounding the omitted tails directly by Gaussian or exponential decay. Integrating over \(b\) in the inner half disk and \(|r|/2<t<|r|/2+1\) also gives a lower bound on \(\phi\) of the form \(c(1+|r|)^{-C}e^{-|r|}\).

The isolated planar operator \(-\Delta_r/2+W\) has ground state \(\varphi\), energy \(-1/2\), and a fixed positive gap above that mode. Indeed the shifted form on test functions \(u\) is \(\int \phi^2 |\nabla(u/\phi)|^2/2\); by closure and local positivity this gives nonnegativity and a simple kernel. A unit sequence perpendicular to the kernel with form tending to zero would be bounded in \(H^1\) and, upon extraction, converge weakly and locally in \(L^2\) (strongly on compact sets). By compact support of \(W\) its limit would be in the kernel, hence zero. But then the \(1/2\) shift of the form precludes convergence of the form values to zero. For the transverse direction \(z\) we use the normalized Gaussian \(g(z)=(\omega/\pi)^{1/4}e^{-\omega z^2/2}\) of \(-\partial_z^2/2+2\pi\rho z^2\) with energy \(\omega/2\), \(\omega=\sqrt{4\pi\rho}\), where \(\rho>0\) is a large fixed integer to be chosen. This mode too has a fixed gap (same positive ground-state identity, with compactness here from confinement).

We will place the planar centers \(u_i\) at scale \(D=\log\lambda\) using the contact layout. For all sufficiently accurate implementations, we require minimum spacing \(\ge0.95D\), nonlink spacing \(\ge1.2D\), and extent \(\operatorname{poly}(n) O(D)\). Define
\[
 \begin{gathered}
 \psi_i(y)=\varphi(r-u_i) g(z),\quad e_*=-1/2+\omega/2,\\
 v_{ij}=\iint\frac{\psi_i(y)^2\psi_j(y')^2}{|y-y'|}\,dy\,dy'
       =v(|u_i-u_j|),\qquad U=v(0).
 \end{gathered}
\]
These constants measure direct electron repulsion: $U$ is the same-site
coefficient and $v_{ij}$ the coefficient between sites $i$ and $j$.
They are bounded uniformly. The function \(v\) is radial as indicated
and uniformly Lipschitz, as follows by translating the smooth density
with decaying derivatives.

To control departures from one electron per site, we need a uniform
lower bound on the Coulomb quadratic form. We prove
\((v_{ij})\ge cI\), with fixed \(c>0\), uniformly over these separated
geometries once the polynomial scale is sufficiently large. For real
coefficients \(a_i\), put \(F=\sum_i a_i\psi_i^2\). Choose a fixed
smooth nonnegative, nonzero test bump near the origin, and let
\(\chi_i\) be its disjoint translates to the centers \((u_i,0)\).
With \(w=\sum_i a_i\chi_i\), Poisson's equation and integration by
parts give
\[
 \left|\int Fw\right|^2\le
 \frac1{4\pi}\left(\iint\frac{F(y)F(y')}{|y-y'|}\,dy\,dy'\right)
       \int|\nabla w|^2
\]
for these smooth decaying densities, or first after truncation at
infinity. To bound the left side below even for signed \(a_i\), let
\(A_{ij}=\int\psi_i^2\chi_j\). Its diagonal entries equal a fixed
\(a_*>0\), while the row and column sums of its off-diagonal entries
are exponentially small in \(D\), up to a polynomial factor in \(m\).
Thus \(A=a_*I+R\), with \(\|R\|\le a_*/2\) at our scales, and
\[
       \int Fw=a^{\mathsf T}Aa\ge(a_*/2)\sum_i a_i^2,
       \qquad \int|\nabla w|^2=O\Bigl(\sum_i a_i^2\Bigr).
\]
The displayed Coulomb inequality proves the required lower bound.
Here and below we take \(\lambda\) so large a power of \(n\) that
exponential-in-\(-D\) bounds beat the polynomial factors. In particular,
applying the bound to two sites and coefficients \((1,-1)\) shows that
\(U-v(d)\) is bounded away from zero for every distance
\(d\ge0.95D\) in use; it is also uniformly bounded above.

\subsection{The occupancy penalty and spin exchange}\label{subsec:finite-fermion}

For the moment keep any geometry satisfying the separation bounds, and
let the hopping strengths \(t_e\) be free parameters. On \(2m\)
canonical fermion modes, two spins at each site, restrict to \(m\)
electrons and write \(n_i\) for total occupancy at site \(i\). For
this finite model we multiply dilated energies by \(\lambda\), so
its Coulomb coefficients have no \(1/\lambda\) factor. The direct
Coulomb interaction is
\[
       \frac U2\sum_i n_i(n_i-1)+\sum_{i<j}v_{ij}n_i n_j.
\]
Set \(s_i=\sum_{j\ne i}v_{ij}\). A diagonal one-body term
\(-\sum_i s_i n_i\) cancels the terms linear in the occupancy defects:
using \(v_{ii}=U\) and \(\sum_i n_i=m\),
\begin{equation}\tag{C}\label{eq:occupancy-completion}
 \begin{aligned}
 &\frac U2\sum_i n_i(n_i-1)+\sum_{i<j}v_{ij}n_i n_j-\sum_i s_i n_i\\
 &\hspace{2em}=\frac12\sum_{i,j}v_{ij}(n_i-1)(n_j-1)-\sum_{i<j}v_{ij}.
 \end{aligned}
\end{equation}
Section~\ref{sec:density} will produce this one-body term by adjusting
the well depths. The finite Hamiltonian to be realized is therefore
\[
 \begin{split}
 H_d(t)={}&\frac12\sum_{i,j}v_{ij}(n_i-1)(n_j-1)-\sum_{i<j}v_{ij}\\
    &-\sum_{e=\{i,j\}}t_e
       \sum_{\sigma=\uparrow,\downarrow}
          (c_{i\sigma}^\dagger c_{j\sigma}+c_{j\sigma}^\dagger c_{i\sigma}).
 \end{split}                                                   \tag{2}\label{eq:2}
\]
The nonnegative quadratic penalty has kernel exactly the singly occupied
spin space. It has a fixed gap: a nonzero integer defect vector
\((n_i-1)_i\) has sum zero, hence squared length at least two, so its
penalty is at least \(c\) by Coulomb coercivity.

Hopping vanishes on first-order compression to single occupancy. Its
second-order effect is the double-occupancy superexchange familiar from
Anderson~\cite[Eqs.~(12)--(13), (18)--(21)]{Anderson1959}; his transfer-energy discussion
also retains intersite Coulomb contributions. The normalization and error
bound below are proved for the precise matrix used here.

One allowed hop on \(e=\{i,j\}\) creates a doubly occupied site and
an empty site. Its defect vector is \((1,-1)\) on that pair and zero
elsewhere, so its excitation energy is exactly \(U-v_{ij}\). Only a
hop on the same link can restore single occupancy in one further step.
For hopping strength \(t_e\), the negative second-order operator is
\[
                  \frac{t_e^2}{U-v_{ij}}(h^s_{ij}-I).
\]
Indeed the triplets do not hop into same-site pairs, while the singlet
hops into two orthogonal doubly occupied states with amplitudes of
magnitude \(\sqrt2\,t_e\). Apply the creation and annihilation
operators to
\((c_{i\uparrow}^\dagger c_{j\downarrow}^\dagger-
c_{i\downarrow}^\dagger c_{j\uparrow}^\dagger)|0\rangle/\sqrt2\)
to obtain these amplitudes. Other sites can be placed after these two in
the ordering: changing the ordering of singly occupied sites gives only
an overall sign on the spin space.

With a common positive parameter \(\tau\) to be fixed below, choose
target hoppings
\[
                  t_e^{\rm tar}=\tau\sqrt{K_e(U-v_{ij})}.
\]
They give the effective spin term
\(\tau^2(H_+-\sum_e K_e)\). Their hopping operator has norm at most
\(\tau\operatorname{poly}(n)\), uniformly over admissible separated
geometries. By Lemma~\ref{lem:perturbation}, the energy remainder is at
most \(C\tau^3\operatorname{poly}(n)\), since the penalty gap is
uniform. Fix a rational inverse-polynomial \(\tau>0\) small enough
that the lemma applies and this remainder is at most
\(\gamma\tau^2/100\). This choice is independent of \(\lambda\)
and of the final centers.

\subsection{Calibrating the separations}

The modes supply the hopping profile
\[
 T(d)=-\int_{\mathbb R^2}\varphi(r)W(r)\varphi(r-d e_1)\,dr
     =\frac{1}{2\|\phi\|_2^2}\int f(r)\phi(r-d e_1)\,dr .
\]
It is positive, uniformly Lipschitz, and lies between
\(c(1+d)^{-C}e^{-d}\) and \(C(1+d)^Ce^{-d}\) by the tail
estimates. The continuum compression will supply hopping
\(t_e=\lambda T(|u_i-u_j|)\) in the units of \(H_d\). Although the
occupancy penalty involves every pair, the target hopping on link
\(e\) depends only on its own distance through \(v_{ij}=v(d)\).
Thus each link requires a scalar solve for
\[
           F_e(d)=\lambda T(d)-\tau\sqrt{K_e(U-v(d))}.
\]
At \(d=(1-\varepsilon/2)D\) and \(d=(1+\varepsilon/2)D\),
respectively, \(F_e\) is positive and negative if the fixed polynomial
exponent in \(\lambda\) is sufficiently large. Indeed the first
term has respective sizes \(\lambda^{\varepsilon/2}\) and
\(\lambda^{-\varepsilon/2}\), up to powers of \(D\), whereas the
target has polynomial upper and inverse-polynomial lower bounds in
\(n\).

Bisection with approximate sign tests finds an inverse-polynomially
small residual: stop if the sign is ambiguous within the evaluation
error, or continue until the bracket width times a Lipschitz bound is
small enough. No monotonicity or derivative lower bound is required.
The link-length construction in Section~\ref{sec:spin}, rescaled by
\(D\), then implements all these distances at once. Its contraction
solves and a final rational approximation give rational centers with the
required separations and hopping accuracy. Section~\ref{sec:implementation}
details the numerical evaluation.

From now on, \(H_d\) means \(H_d(t)\) at these final centers with
\(t_e=\lambda T(|u_i-u_j|)\). Choose the per-link calibration accuracy
so that replacing \(t_e^{\rm tar}\) by \(t_e\) changes the total
Hamiltonian norm by at most \(\gamma\tau^2/100\). Together with the
second-order remainder, this gives
\[
  |E(H_d)+\sum_{i<j}v_{ij}
       -\tau^2(E(H_+)-\sum_e K_e)|\le 2\gamma\tau^2/100.       \tag{3}\label{eq:3}
\]

\subsection{Polynomial scale convention}\label{subsec:scales}

Equation~\eqref{eq:3} fixes the finite-model accuracy required of the
continuum construction. All subsequent errors must fit inside an
inverse-polynomial fraction of $\gamma\tau^2$ in the energy units of
$H_d$.

In the estimates below \(p(n,D)\) denotes a polynomial upper bound with uniform coefficients and exponents, which may be enlarged between uses. All requirements of the form "make \(p(n,D)\lambda^{-\alpha}\) small" for fixed \(\alpha>0\), with required smallness at worst inverse polynomial in \(n\), can be imposed simultaneously. Specifically we take an integer \(k\) growing as a sufficiently large fixed power of \(n\) (allowing a fixed factor) and put
\[
                  h=1/k,\qquad \lambda=8k^3/\rho .
\]
Here \(h\) will be the side length of the cubature cells in
Section~\ref{sec:nuclei}; the relation
\(\rho h^3/8=1/\lambda\) gives each of their eight nodes the mass
needed for a unit nucleus. Other magnitudes such as the large box used
below will have stated polynomial bounds in \(\lambda,n\); its volume
will enter \(p(n,D)\) only when explicitly accounted for. The fixed
\(\rho\) only has to dominate bounded potential derivatives as described
next and can be chosen before these scale choices.
